\section{Global ideal sums in the theta estimates}
\label{gt-section}

The local results in the preceding sections leave several global
estimates to be proved.  We begin with a complete reconstruction of the
sum over the dual lattices in the scalar orbital calculation.  This
part admits a shorter proof and a stronger estimate than the one
stated in the source~\cite{alpo-uniform-source}.  All the constants in this section are explicit
and independent of the number field and its discriminant.

\subsection{Ideal norms and separation of algebraic integers}

Let $E$ be a totally complex number field of degree $2f$, where $f\geq1$.
Write $\mathcal O_E$ for its ring of integers and choose embeddings
$\sigma_1,\ldots,\sigma_f:E\longrightarrow\mathbb C$, one from each
complex conjugate pair.  A nonzero fractional ideal $\mathfrak a$ is a
nonzero finitely generated $\mathcal O_E$-submodule of $E$ for which
$m\mathfrak a\subseteq\mathcal O_E$ for some positive integer $m$.
Its absolute norm is
\[
 N\mathfrak a=\frac{[\mathcal O_E:m\mathfrak a]}{m^{2f}}.
\]
This does not depend on $m$.  Equivalently it is the ratio of the
covolumes of $\mathfrak a$ and $\mathcal O_E$ under the embedding
$\alpha\mapsto(\sigma_v\alpha)_{v=1}^f$ into $\mathbb C^f$.
The choice of Lebesgue measure cancels in this ratio.

We use the usual multiplicativity of ideal norms.  To recall its
content, every nonzero fractional ideal has a unique factorization
$\prod_{\mathfrak p}\mathfrak p^{n_{\mathfrak p}}$, with finitely many
nonzero integers $n_{\mathfrak p}$, and its norm is
$\prod_{\mathfrak p}|\mathcal O_E/\mathfrak p|^{n_{\mathfrak p}}$.
For integral ideals this follows by the Chinese remainder theorem
and the filtration by successive powers of each prime ideal; each
successive quotient has the cardinality of its residue field.
Clearing denominators gives the fractional-ideal assertion.

\begin{lemma}\label{gt-product-minimum}
If $0\ne\alpha\in\mathfrak a$, then
\begin{equation}\label{gt-product-minimum-eq}
 \prod_{v=1}^f|\sigma_v\alpha|\geq\sqrt{N\mathfrak a}.
\end{equation}
The same inequality applies to $\alpha-\beta$ for distinct
$\alpha,\beta\in\mathfrak a$.
\end{lemma}

\begin{proof}
The containment $\alpha\mathcal O_E\subseteq\mathfrak a$ has a
positive integral index.  Multiplication by $\alpha$ has real
determinant
$|N_{E/\mathbb Q}\alpha|=\prod_v|\sigma_v\alpha|^2$ on
$\mathbb C^f$.  Comparing covolumes gives
\[
 |N_{E/\mathbb Q}\alpha|
   =N\mathfrak a\,[\mathfrak a:\alpha\mathcal O_E]
   \geq N\mathfrak a.
\]
Taking square roots proves the first assertion.  The second follows
because $\mathfrak a$ is an additive subgroup.
\end{proof}

\begin{lemma}[A box count]\label{gt-box-count}
For positive real numbers $R_1,\ldots,R_f$, let
\[
 B_{\mathfrak a}(R)=
 \#\{0\ne\alpha\in\mathfrak a:
                  |\sigma_v\alpha|\leq R_v\text{ for all }v\}.
\]
Then
\begin{equation}\label{gt-box-count-eq}
 B_{\mathfrak a}(R)
       \leq9^f\frac{\prod_vR_v^2}{N\mathfrak a}.
\end{equation}
Moreover $B_{\mathfrak a}(R)=0$ if
$\prod_vR_v<\sqrt{N\mathfrak a}$.
\end{lemma}

\begin{proof}
The last assertion is Lemma~\ref{gt-product-minimum}.
If the set is empty, the estimate is immediate.  Otherwise put
\[
 a=\left(\frac{\sqrt{N\mathfrak a}}{\prod_vR_v}\right)^{1/f},
 \qquad 0<a\leq1.
\]
After dividing the $v$th coordinate by $R_v$, the points lie in the
closed unit polydisc in $\mathbb C^f$.  Distinct resulting points
$z,z'$ satisfy
\[
 \prod_v|z_v-z'_v|\geq a^f,
 \qquad \max_v|z_v-z'_v|\geq a.
\]
The open polydiscs of radius $a/2$ about these points are therefore
pairwise disjoint.  They are all contained in the polydisc of radius
$1+a/2$.  Comparing their $2f$-dimensional volumes gives
\[
 B_{\mathfrak a}(R)
 \leq\left(1+\frac2a\right)^{2f}
 \leq\left(\frac3a\right)^{2f}
 =9^f\frac{\prod_vR_v^2}{N\mathfrak a}.
\]
One may apply this argument first to every finite subset.  The
uniform finite bound then also proves finiteness of the full set,
so no separate lattice-point asymptotic is being used.
\end{proof}

\subsection{A stronger smoothed sum with its exact logarithmic factor}

For $x=(x_v)_{v=1}^f$ with $x_v>0$, and $0\ne\alpha\in E$, set
\[
 q_x(\alpha)=\prod_{v=1}^f
       \min\left(1,\frac{x_v}{|\sigma_v\alpha|}\right),
 \qquad
 W_{K,x}(\alpha)=q_x(\alpha)^K
 =\prod_{v=1}^f
       \min\left(1,\left(\frac{x_v}{|\sigma_v\alpha|}\right)^K\right).
\]
Here $K$ is a positive real number.  Define $\Theta>0$ by
\begin{equation}\label{gt-theta}
 \Theta^f=\frac{\sqrt{N\mathfrak a}}{\prod_vx_v}.
\end{equation}

\begin{theorem}\label{gt-smoothed}
If $K>2$ and $\Theta\geq1$, then the sum over all nonzero elements
of $\mathfrak a$ converges and satisfies
\begin{equation}\label{gt-smoothed-exact}
 \sum_{0\ne\alpha\in\mathfrak a}W_{K,x}(\alpha)
 \leq
 \left(\frac{9K}{K-2}\right)^f\Theta^{-Kf}
 \sum_{j=0}^{f-1}\frac{((K-2)f\log\Theta)^j}{j!}.
\end{equation}
For $K\geq3$ and $\Theta\geq1$, the following simpler estimates
also hold:
\begin{align}
 \sum_{0\ne\alpha\in\mathfrak a}W_{K,x}(\alpha)
    &\leq\bigl(27\Theta^{-(K-1)}\bigr)^f,
       \label{gt-smoothed-simple}\\
 \sum_{0\ne\alpha\in\mathfrak a}W_{K,x}(\alpha)
    &\leq\bigl(81(1+\log\Theta)\Theta^{-K}\bigr)^f.
       \label{gt-smoothed-log}
\end{align}
\end{theorem}

\begin{proof}
For $r>0$ and $x>0$ the elementary identity
\[
 \min(1,x/r)^K
       =K\int_1^\infty 1_{\{r\leq xt\}}t^{-K-1}\,dt
\]
follows by integrating from $\max(1,r/x)$ to infinity.
Multiply this identity over the embeddings and sum over the ideal.
All functions are nonnegative, so the interchange of sums and
integrals is justified by monotone convergence.  The box count gives
\begin{align*}
 \sum_{0\ne\alpha\in\mathfrak a}W_{K,x}(\alpha)
 &=K^f\int_{[1,\infty)^f}
        B_{\mathfrak a}((x_vt_v)_v)
                 \prod_vt_v^{-K-1}\,dt_1\cdots dt_f\\
 &\leq9^fK^f\Theta^{-2f}
       \int_{\substack{t_v\geq1\\\prod_vt_v\geq\Theta^f}}
                         \prod_vt_v^{-(K-1)}\,dt_1\cdots dt_f.
\end{align*}
The restriction on the product is justified by the vanishing
assertion of Lemma~\ref{gt-box-count}.  Substitute $t_v=e^{u_v}$.
Writing $a=K-2>0$ and $L=f\log\Theta\geq0$, the last integral is
\[
 \int_{\substack{u_v\geq0\\\sum_vu_v\geq L}}
                      e^{-a\sum_vu_v}\,du_1\cdots du_f
   =\frac1{(f-1)!}\int_L^\infty e^{-au}u^{f-1}\,du.
\]
The equality follows by first computing the volume
$u^f/f!$ of $\{u_v\geq0:\sum_vu_v\leq u\}$ and then integrating
with respect to $u$.  Repeated integration by parts gives
\[
 \frac1{(f-1)!}\int_L^\infty e^{-au}u^{f-1}\,du
       =a^{-f}e^{-aL}\sum_{j=0}^{f-1}\frac{(aL)^j}{j!}.
\]
Substituting $e^{-aL}=\Theta^{-(K-2)f}$ proves
\eqref{gt-smoothed-exact}, including convergence.  These formulas
also hold for $f=1$, when the polynomial is $1$.

For the two simpler estimates, the same integral calculation without
the product restriction shows, for every $s>2$,
\begin{equation}\label{gt-moment-bound}
 \sum_{0\ne\alpha\in\mathfrak a}q_x(\alpha)^s
       \leq\left(\frac{9s}{s-2}\right)^f\Theta^{-2f}.
\end{equation}
Lemma~\ref{gt-product-minimum} also gives
\begin{equation}\label{gt-q-bound}
 q_x(\alpha)\leq
       \frac{\prod_vx_v}{\prod_v|\sigma_v\alpha|}
       \leq\Theta^{-f}.
\end{equation}
Consequently, if $s\leq K$, multiplication of
\eqref{gt-moment-bound} by $\Theta^{-f(K-s)}$ yields
\begin{equation}\label{gt-interpolation}
 \sum_{0\ne\alpha\in\mathfrak a}q_x(\alpha)^K
       \leq\left(\frac{9s}{s-2}
                         \Theta^{-(K-s+2)}\right)^f.
\end{equation}
Taking $s=3$ proves \eqref{gt-smoothed-simple}.
For \eqref{gt-smoothed-log}, take
$s=2+\eta$, where $\eta=(1+\log\Theta)^{-1}$.
Then $0<\eta\leq1$, so $s\leq3\leq K$, and
\[
 \frac{9s}{s-2}=9(3+2\log\Theta)
       \leq27(1+\log\Theta),\qquad
 \Theta^\eta\leq e<3.
\]
Equation~\eqref{gt-interpolation} now gives the stated bound.
\end{proof}

\begin{corollary}\label{gt-subgroup-sum}
Let $L\subseteq\mathfrak a$ be an additive subgroup, let $A>0$,
and suppose that $x_v\geq0$ and
\begin{equation}\label{gt-product-size}
 \prod_vx_v\leq(A/\Theta)^f\sqrt{N\mathfrak a},
 \qquad \Theta\geq A.
\end{equation}
For $K>2$, the right side of \eqref{gt-smoothed-exact}, with
$\Theta$ replaced by $\Theta/A$, bounds the sum of $W_{K,x}$ over
$L\setminus\{0\}$.  For $K\geq3$ one may also use
\begin{equation}\label{gt-subgroup-simple}
 \bigl(27(A/\Theta)^{K-1}\bigr)^f
 \quad\hbox{or}\quad
 \bigl(81(1+\log(\Theta/A))(A/\Theta)^K\bigr)^f.
\end{equation}
Every assertion holds, in particular, for every finite partial sum.
\end{corollary}

\begin{proof}
If some $x_v=0$, every summand is zero: a nonzero field element has
nonzero image under every embedding.  Suppose therefore that all
$x_v>0$.  Define $T$ by
$T^f=\sqrt{N\mathfrak a}/\prod_vx_v$.  Then
$T\geq\Theta/A\geq1$, and Theorem~\ref{gt-smoothed} applies.
The exact expression in \eqref{gt-smoothed-exact} is decreasing in
$T$: before its evaluation it is the product of $T^{-2f}$ and a
positive integral over
$\{u_v\geq0:\sum_vu_v\geq f\log T\}$.
The other two expressions are also decreasing for $T\geq1$ and
$K\geq3$.  For the logarithmic expression this follows directly from
\[
 \frac{d}{dT}\bigl((1+\log T)T^{-K}\bigr)
       =T^{-K-1}\bigl(1-K(1+\log T)\bigr)<0.
\]
The sum over $L$ is at most the sum over $\mathfrak a$.
\end{proof}

\subsection{The actual period and dual lattices}

We now assume that $E$ is a CM field and write
$\alpha\mapsto\bar\alpha$ for its complex conjugation.  Put
\[
 \Delta_E=\operatorname{disc}(E),\qquad
 \rho_E=|\Delta_E|^{1/(2f)},\qquad
 \mathfrak D_E^{-1}
   =\{x\in E:\operatorname{Tr}_{E/\mathbb Q}
                          (x\mathcal O_E)\subseteq\mathbb Z\}.
\]
Thus $\mathfrak D_E$ is the different, and $\rho_E$ is the root
discriminant.  The trace form is nondegenerate in characteristic
zero, so the displayed set is a nonzero fractional ideal.

\begin{lemma}\label{gt-different}
For every nonzero fractional ideal $\mathfrak a$,
\begin{align}
 \{x\in E:\operatorname{Tr}_{E/\mathbb Q}(x\mathfrak a)
                    \subseteq\mathbb Z\}
       &=\mathfrak D_E^{-1}\mathfrak a^{-1},
          \label{gt-trace-dual}\\
 N\mathfrak D_E&=|\Delta_E|.\label{gt-different-norm}
\end{align}
The ideals $\mathfrak D_E^\nu$ are preserved by complex
conjugation for every integer $\nu$.
\end{lemma}

\begin{proof}
If $x$ belongs to the set on the left of \eqref{gt-trace-dual},
then for every $y\in\mathfrak a$ and $b\in\mathcal O_E$ one has
$by\in\mathfrak a$ and $\operatorname{Tr}(xyb)\in\mathbb Z$.
Thus $x\mathfrak a\subseteq\mathfrak D_E^{-1}$.  Conversely
this containment implies the trace condition by taking $b=1$.
Multiplication by the inverse ideal gives
\eqref{gt-trace-dual}.

Let $e_1,\ldots,e_{2f}$ be an integral basis of $\mathcal O_E$
and let $T=(\operatorname{Tr}(e_ie_j))_{i,j}$.  In these coordinates
the codifferent is $T^{-1}\mathbb Z^{2f}$.  Its covolume relative
to $\mathcal O_E$ is $|\det T|^{-1}=|\Delta_E|^{-1}$.
Multiplicativity of norms gives \eqref{gt-different-norm}.
Finally the trace is unchanged by complex conjugation.  Hence
conjugation preserves the codifferent, its inverse, and all their
integer powers.
\end{proof}

Fix a positive integer $g$, an integer $\nu$, and a nonnegative
integer $s$.  The period lattice and the lift lattice in the global
orbital calculation are, respectively,
\begin{align}
 P_{\nu,s}
   &=\{y\in g\mathfrak D_E^\nu:
                    y-(-1)^s\bar y\in g^2\mathfrak D_E^\nu\},
           \label{gt-period-lattice}\\
 L_{\nu,s}
   &=\{x\in g^{-1}\mathfrak D_E^{-\nu}:
                    x-(-1)^s\bar x\in\mathfrak D_E^{-\nu}\}.
           \label{gt-lift-lattice}
\end{align}
The dual of the period lattice is taken for the pairing
$\operatorname{Tr}_{E/\mathbb Q}(xy/g)$:
\begin{equation}\label{gt-period-dual-definition}
 P_{\nu,s}^{\vee,g}
   =\{x\in E:\operatorname{Tr}_{E/\mathbb Q}(xy/g)
                    \in\mathbb Z\text{ for every }y\in P_{\nu,s}\}.
\end{equation}
The superscripts specify this normalization; using the unscaled
trace pairing would give a different power of $g$.

\begin{proposition}\label{gt-lattice-containments}
For the lattices just defined,
\begin{align}
 g^2\mathfrak D_E^\nu
       &\subseteq P_{\nu,s}\subseteq g\mathfrak D_E^\nu,
             \label{gt-period-between}\\
 \mathfrak D_E^{-\nu-1}
       &\subseteq P_{\nu,s}^{\vee,g}
                  \subseteq g^{-1}\mathfrak D_E^{-\nu-1},
             \label{gt-dual-between}\\
 \mathfrak D_E^{-\nu}
       &\subseteq L_{\nu,s}
                  \subseteq g^{-1}\mathfrak D_E^{-\nu}.
             \label{gt-lift-between}
\end{align}
In particular these groups are full lattices.  The containing ideals
in \eqref{gt-dual-between} and \eqref{gt-lift-between} have norms
\begin{align}
 \sqrt{N(g^{-1}\mathfrak D_E^{-\nu-1})}
       &=(g\rho_E^{\nu+1})^{-f},
             \label{gt-dual-norm}\\
 \sqrt{N(g^{-1}\mathfrak D_E^{-\nu})}
       &=(g\rho_E^\nu)^{-f}.
             \label{gt-lift-norm}
\end{align}
These assertions need no unramifiedness assumption on $E$ and no
restriction on the primes dividing $2g$.
\end{proposition}

\begin{proof}
The different and all its powers are conjugation invariant.  Thus
$y\in g^2\mathfrak D_E^\nu$ implies both
$y\in g\mathfrak D_E^\nu$ and
$y-(-1)^s\bar y\in g^2\mathfrak D_E^\nu$.
This proves \eqref{gt-period-between}.
If $x\in P_{\nu,s}^{\vee,g}$, its trace pairing with
$g^2\mathfrak D_E^\nu$ is integral, so
\[
 \operatorname{Tr}_{E/\mathbb Q}(gx\mathfrak D_E^\nu)
        \subseteq\mathbb Z.
\]
Lemma~\ref{gt-different} gives
$gx\in\mathfrak D_E^{-\nu-1}$, proving the upper containment in
\eqref{gt-dual-between}.  For its lower containment, pair
$x\in\mathfrak D_E^{-\nu-1}$ with
$y\in P_{\nu,s}\subseteq g\mathfrak D_E^\nu$.
Then $xy/g\in\mathfrak D_E^{-1}$, whose trace is integral.
The upper containment in \eqref{gt-lift-between} is part of the
definition, and the lower one follows from conjugation invariance
in the same way as \eqref{gt-period-between}.

The lower and upper groups in these containments are nonzero
fractional ideals, so each intervening additive group is a full
lattice.  Finally $N(g\mathcal O_E)=g^{2f}$ and
$N\mathfrak D_E=\rho_E^{2f}$.  Multiplicativity proves
\eqref{gt-dual-norm} and \eqref{gt-lift-norm} for every integer
$\nu$, including negative $\nu$.
\end{proof}

For the blocks used in the theta calculation, the exponent is
\begin{equation}\label{gt-block-exponent}
 \nu(r,b,s)=2r+3-2b-s.
\end{equation}
Here $r$ is a nonnegative integer, $b$ is the first block index, and
$s\geq1$ is its length.  Define
\[
 \mathfrak a_{b,s}=g^{-1}\mathfrak D_E^{-(\nu(r,b,s)+1)},
 \qquad
 \mathfrak l_{b,s}=g^{-1}\mathfrak D_E^{-\nu(r,b,s)}.
\]
Then
\begin{equation}\label{gt-adjacent-ideals}
 \mathfrak l_{b,s}=\mathfrak a_{b,s+1},\qquad
 \sqrt{N\mathfrak a_{b,s}}
    =(g\rho_E^{2r+4-2b-s})^{-f},\qquad
 \sqrt{N\mathfrak l_{b,s}}
    =(g\rho_E^{2r+3-2b-s})^{-f}.
\end{equation}
This is an identity of ideals, not an estimate; it accounts for the
one-step change in the discriminant exponent in the recursion.

\begin{corollary}[The actual dual-lattice error]\label{gt-actual-error}
Let $\Theta>4$.  Suppose that the nonnegative widths
$x^{D}_{b,v}$ and $x^{L}_{b,s,v}$ satisfy
\begin{align}
 \prod_vx^{D}_{b,v}
       &\leq\Theta^{-f}(g\rho_E^{2r+3-2b})^{-f},
                 \label{gt-dual-width}\\
 \prod_vx^{L}_{b,s,v}
       &\leq(2/\Theta)^f
                       (g\rho_E^{2r+3-2b-s})^{-f}.
                 \label{gt-lift-width}
\end{align}
For each indicated block, the sums
\begin{equation}\label{gt-actual-sums}
 \sum_{0\ne\alpha\in P_{\nu(r,b,1),1}^{\vee,g}}
                  W_{5,x^D_b}(\alpha),
 \qquad
 \sum_{0\ne\alpha\in L_{\nu(r,b,s),s}}
                  W_{5,x^L_{b,s}}(\alpha)
\end{equation}
are bounded by each of the following quantities:
\begin{align}
 &15^f(2/\Theta)^{5f}
       \sum_{j=0}^{f-1}
                  \frac{(3f\log(\Theta/2))^j}{j!},
                  \label{gt-error-exact}\\
 &(432\Theta^{-4})^f,
                  \label{gt-error-simple}\\
 &\bigl(2592(1+\log(\Theta/2))\Theta^{-5}\bigr)^f.
                  \label{gt-error-log}
\end{align}
The bounds are uniform over the block indices and over every finite
family of width vectors satisfying
\eqref{gt-dual-width}--\eqref{gt-lift-width}.
\end{corollary}

\begin{proof}
Proposition~\ref{gt-lattice-containments} places the two summation
groups in $\mathfrak a_{b,1}$ and $\mathfrak l_{b,s}$,
respectively.  Their norm formulas are
\eqref{gt-adjacent-ideals}.  Thus the assumptions are precisely
\eqref{gt-product-size}, with $A=1$ for the first group and $A=2$
for the second.  Corollary~\ref{gt-subgroup-sum}, with $K=5$, gives
the asserted common bounds by using $A=2$ for both.  The constants
are $9K/(K-2)=15$, $27\cdot2^4=432$, and
$81\cdot2^5=2592$.  No union bound over the blocks is taken:
the same bound holds separately for every required sum.
\end{proof}

\subsection{Exact covolumes, including primes above two}

Write $F$ for the maximal totally real subfield of $E$, so
$[F:\mathbb Q]=f$.  In this subsection suppose that $E/F$ is
unramified at every finite prime of $F$ dividing $g$.
This is weaker than assuming that $E/F$ is unramified at every
finite prime.  It also holds when every rational prime dividing
$g$ splits completely in $E$, as at the selected levels in the
application.  We give the elementary local calculation, since
division by $2$ would be invalid at some of these primes.

\begin{lemma}\label{gt-quadratic-count}
Let $A$ be the ring of integers of a nonarchimedean local field,
let $B$ be a quadratic \emph{\'etale} $A$-algebra which is
unramified, and let $\tau$ be its nontrivial involution.
Thus $B=A\times A$, with the factors exchanged, or $B$ is the
ring of integers of the unramified quadratic field extension.
For a nonzero nonunit $a\in A$ and $\epsilon\in\{1,-1\}$,
\begin{equation}\label{gt-local-kernel-card}
 \#\ker\bigl(1-\epsilon\tau:B/aB\longrightarrow B/aB\bigr)
       =\#(A/aA).
\end{equation}
\end{lemma}

\begin{proof}
In the split case the kernel consists of the pairs
$(x,\epsilon x)$, so the assertion is immediate.
In the field case choose an element of the quadratic residue
field which generates it over the residue field of $A$, and lift
it to $\omega\in B$.  Nakayama's lemma shows that
$1,\omega$ is an $A$-basis of $B$.  Its conjugate has a different
residue, since finite fields are perfect, so
$\delta=\omega-\tau\omega$ is a unit of $B$.
Put $t=\omega+\tau\omega\in A$.

For $\epsilon=1$, write an element as $x+y\omega$.
Its difference with its conjugate is $y\delta$, which vanishes
modulo $aB$ if and only if $y\in aA$.  The kernel is therefore
the copy of $A/aA$ represented by $x$.
For $\epsilon=-1$ the map is the trace
$x+y\omega\mapsto2x+ty$ into $A/aA$.
This map is surjective: if the residue characteristic is odd,
$2$ is a unit; if it is $2$, the equality
$\delta=2\omega-t$ shows that $t$ is a unit.
Its domain has cardinality $\#(A/aA)^2$, so its kernel has
cardinality $\#(A/aA)$.  This proves both signs without
inverting $2$.
\end{proof}

\begin{proposition}\label{gt-exact-covolumes}
Under the unramifiedness assumption of this subsection,
\begin{align}
 [g\mathfrak D_E^\nu:P_{\nu,s}]
    &=[P_{\nu,s}:g^2\mathfrak D_E^\nu]=g^f,
                 \label{gt-period-indices}\\
 [g^{-1}\mathfrak D_E^{-\nu}:L_{\nu,s}]&=g^f.
                 \label{gt-lift-index}
\end{align}
For standard Lebesgue measure on $\mathbb C^f$, the covolumes
under $\sigma=(\sigma_v)_v$ are
\begin{align}
 \operatorname{covol}\sigma(P_{\nu,s})
    &=2^{-f}g^{3f}\rho_E^{(2\nu+1)f},
                 \label{gt-period-covol}\\
 \operatorname{covol}\sigma(L_{\nu,s})
    &=2^{-f}g^{-f}\rho_E^{(1-2\nu)f}.
                 \label{gt-lift-covol}
\end{align}
In particular, if $\Lambda_{\nu,s}=\rho_E^{-\nu}
\sigma(P_{\nu,s})$, then
\begin{equation}\label{gt-normalized-covol}
 \operatorname{covol}\Lambda_{\nu,s}
      =(\rho_E/2)^f g^{3f},
\end{equation}
independently of $\nu$ and $s$.
\end{proposition}

\begin{proof}
At a finite prime $\mathfrak p$ of $F$ dividing $g$, put
$A=\mathcal O_{F,\mathfrak p}^{\wedge}$ and
$B=\mathcal O_E\otimes_{\mathcal O_F}A$.
The trace pairing of $B$ over $A$ is perfect because $B/A$ is
unramified.  This can be seen directly from the basis in
Lemma~\ref{gt-quadratic-count}: its trace determinant is
$(\omega-\tau\omega)^2$, a unit; in the split case use the
two standard idempotents.  The absolute trace pairing is the
composition of this relative pairing with the trace of $F$.
Consequently the localization of $\mathfrak D_E$ is
$\mathfrak D_F B$.  More explicitly, choose a basis of $B$
over $A$ and its relative trace-dual basis, both integral.
An element is dual to $B$ for the absolute trace exactly when
all its coefficients in that dual basis belong to the absolute
codifferent of $A$.  This proves the assertion about the
different and its inverse, and hence about every integer power.
Thus locally $\mathfrak D_E^\nu=dB$ for a nonzero $d\in F_{\mathfrak p}$
fixed by $\tau$.

For $\epsilon=(-1)^s$, reduction modulo $g\mathfrak D_E^\nu$
therefore identifies
\[
 \{z\in\mathfrak D_E^\nu:z-\epsilon\bar z
                            \in g\mathfrak D_E^\nu\}
       /g\mathfrak D_E^\nu
\]
with a product, over $\mathfrak p\mid g$, of the kernels in
Lemma~\ref{gt-quadratic-count}.  Its cardinality is
$\#(\mathcal O_F/g\mathcal O_F)=g^f$.
The whole quotient
$\mathfrak D_E^\nu/g\mathfrak D_E^\nu$ has cardinality
$g^{2f}$.  Multiplication by $g$ gives both indices in
\eqref{gt-period-indices}.  Applying the same argument with
$\mathfrak D_E^{-\nu}$ and then multiplying by $g^{-1}$ proves
\eqref{gt-lift-index}.  When $g=1$ the quotients are trivial
and the same formulas hold.

For completeness,
\begin{equation}\label{gt-ideal-covol}
 \operatorname{covol}\sigma(\mathfrak a)
       =2^{-f}|\Delta_E|^{1/2}N\mathfrak a
       =2^{-f}\rho_E^fN\mathfrak a.
\end{equation}
Indeed the matrix of all $2f$ complex embeddings on an integral
basis has determinant of absolute value $|\Delta_E|^{1/2}$.
Replacing each conjugate pair of rows by its real and imaginary
parts divides the determinant's absolute value by $2$.
This gives the formula for $\mathcal O_E$, and the definition of
$N\mathfrak a$ gives it for every fractional ideal.
The two index formulas and
$N(g^a\mathfrak D_E^b)=g^{2af}\rho_E^{2bf}$ now give
\eqref{gt-period-covol} and \eqref{gt-lift-covol}.
Multiplication by the positive real number $\rho_E^{-\nu}$
on $\mathbb C^f$ scales volume by $\rho_E^{-2\nu f}$,
which proves \eqref{gt-normalized-covol}.
\end{proof}

Use the real inner product
$\langle z,w\rangle_{\mathbb R}
   =\operatorname{Re}\sum_v\bar z_vw_v$ on $\mathbb C^f$.
For a full lattice $\Lambda$, write
\[
 \Lambda^*=\{\xi\in\mathbb C^f:
       \langle\xi,\lambda\rangle_{\mathbb R}\in\mathbb Z
                       \text{ for every }\lambda\in\Lambda\}.
\]

\begin{proposition}\label{gt-Fourier-dual}
For every nonzero real number $c$, without any unramifiedness
assumption, one has
\begin{equation}\label{gt-Fourier-dual-eq}
 (c\sigma(P_{\nu,s}))^*
       =\left\{\frac2{gc}\overline{\sigma(x)}:
                             x\in P_{\nu,s}^{\vee,g}\right\}.
\end{equation}
Under the unramifiedness assumption used in
Proposition~\ref{gt-exact-covolumes}, one further has
\begin{align}
 P_{\nu,s}^{\vee,g}&=L_{\nu+1,s+1},
                 \label{gt-exact-trace-dual}\\
 \operatorname{covol}\sigma(P_{\nu,s}^{\vee,g})
       &=2^{-f}g^{-f}\rho_E^{-(2\nu+1)f},
                 \label{gt-dual-covol}\\
 \operatorname{covol}\bigl(\rho_E^{\nu+1}
                       \sigma(P_{\nu,s}^{\vee,g})\bigr)
       &=(\rho_E/(2g))^f.
                 \label{gt-normalized-dual-covol}
\end{align}
\end{proposition}

\begin{proof}
For $x,y\in E$ the normalizations satisfy exactly
\begin{equation}\label{gt-pairing-identity}
 \left\langle\frac2{gc}\overline{\sigma(x)},
                         c\sigma(y)\right\rangle_{\mathbb R}
       =\frac2g\operatorname{Re}\sum_v\sigma_v(xy)
       =\operatorname{Tr}_{E/\mathbb Q}(xy/g).
\end{equation}
This gives one inclusion in \eqref{gt-Fourier-dual-eq}.
For the other, choose a $\mathbb Z$-basis $y_1,\ldots,y_{2f}$
of $P_{\nu,s}$.  It is a $\mathbb Q$-basis of $E$.
If $\xi$ belongs to the Fourier dual, the prescribed values
$\langle\xi,c\sigma(y_i)\rangle_{\mathbb R}$ are integers.
Nondegeneracy of the rational trace form gives a unique $x\in E$
with those values equal to $\operatorname{Tr}(xy_i/g)$.
The $2f$ vectors $c\sigma(y_i)$ form a real basis, so
\eqref{gt-pairing-identity} implies
$\xi=2\overline{\sigma(x)}/(gc)$.
The integral values show that $x\in P_{\nu,s}^{\vee,g}$.

For any real full lattice, its Euclidean dual has reciprocal
covolume.  In coordinates, if the basis matrix is $M$, the dual
basis matrix is $(M^{-1})^{\mathsf T}$, which proves this assertion.
Taking $c=1$ in \eqref{gt-Fourier-dual-eq} and using
\eqref{gt-period-covol} therefore gives
\[
 (2/g)^{2f}\operatorname{covol}\sigma(P_{\nu,s}^{\vee,g})
       =\bigl(2^{-f}g^{3f}\rho_E^{(2\nu+1)f}\bigr)^{-1}.
\]
This is \eqref{gt-dual-covol}.

It remains to identify the lattice, rather than just its volume.
Put $\epsilon=(-1)^s$.  If $x\in P_{\nu,s}^{\vee,g}$, then
$x\in g^{-1}\mathfrak D_E^{-\nu-1}$ by
Proposition~\ref{gt-lattice-containments}.
For every $z\in\mathfrak D_E^\nu$ the element
$g(z+\epsilon\bar z)$ belongs to $P_{\nu,s}$.
Pairing with it and using invariance of trace under conjugation
gives
\[
 \operatorname{Tr}\bigl((x+\epsilon\bar x)z\bigr)\in\mathbb Z.
\]
Hence $x+\epsilon\bar x\in\mathfrak D_E^{-\nu-1}$ by
Lemma~\ref{gt-different}, and so
$P_{\nu,s}^{\vee,g}\subseteq L_{\nu+1,s+1}$.
Equation~\eqref{gt-lift-covol} shows that the latter lattice has
the covolume in \eqref{gt-dual-covol}.  The inclusion has index
one, proving \eqref{gt-exact-trace-dual}.  Finally scaling by
$\rho_E^{\nu+1}$ gives \eqref{gt-normalized-dual-covol}.
\end{proof}

\subsection{The global Poisson inequality with its actual regularity}

The summation step in the source uses continuous nonnegative
functions, not merely Schwartz functions.  We prove the needed
version.  This also fixes the Fourier and volume normalizations.
For an integrable function $\phi:\mathbb C^f\to\mathbb C$, let
\[
 \widehat\phi(\xi)=\int_{\mathbb C^f}\phi(z)
          e^{-2\pi i\langle\xi,z\rangle_{\mathbb R}}\,dz,
\]
where $dz$ is standard Lebesgue measure.

\begin{theorem}\label{gt-Poisson}
Let $\Lambda$ be a full lattice in $\mathbb C^f$ and
$V=\operatorname{covol}\Lambda$.
Suppose that $\phi:\mathbb C^f\to\mathbb C$ is continuous and
that, for some finite $C\geq0$,
\begin{equation}\label{gt-Poisson-decay}
 |\phi(z)|\leq C\prod_{v=1}^f(1+|z_v|)^{-3}.
\end{equation}
If $\sum_{\xi\in\Lambda^*}|\widehat\phi(\xi)|<\infty$, then
for every $z_0\in\mathbb C^f$,
\begin{equation}\label{gt-Poisson-eq}
 \sum_{\lambda\in\Lambda}\phi(z_0+\lambda)
       =V^{-1}\sum_{\xi\in\Lambda^*}
          \widehat\phi(\xi)
                 e^{2\pi i\langle\xi,z_0\rangle_{\mathbb R}}.
\end{equation}
Both series converge absolutely, and the first converges locally
uniformly in $z_0$.
If $\phi$ is real and nonnegative, then every finite subset
$S\subseteq z_0+\Lambda$ satisfies
\begin{equation}\label{gt-Poisson-inequality}
 \sum_{z\in S}\phi(z)
       \leq V^{-1}\int_{\mathbb C^f}\phi(z)\,dz
          +V^{-1}\sum_{0\ne\xi\in\Lambda^*}
                                  |\widehat\phi(\xi)|.
\end{equation}
\end{theorem}

\begin{proof}
Write $h(z)=\prod_v(1+|z_v|)^{-3}$.
This function is integrable because each complex coordinate has
real dimension two and
$2\pi\int_0^\infty r(1+r)^{-3}\,dr<\infty$.
Choose $\delta>0$ so that the open polydiscs
$\lambda+\{z:\max_v|z_v|<\delta\}$, $\lambda\in\Lambda$,
are pairwise disjoint.  Such a choice exists because a full
lattice is discrete.  On each of these polydiscs,
\[
 h(\lambda)\leq(1+\delta)^{3f}h(z).
\]
Integration over the disjoint polydiscs proves
$\sum_{\lambda\in\Lambda}h(\lambda)<\infty$.
For $z_0$ in any fixed compact set, a further fixed factor bounds
$h(z_0+\lambda)$ by $h(\lambda)$, by the triangle inequality.
Thus \eqref{gt-Poisson-decay} and the Weierstrass test imply that
\[
 F(z_0)=\sum_{\lambda\in\Lambda}\phi(z_0+\lambda)
\]
is a continuous $\Lambda$-periodic function.

Integrate over a half-open fundamental parallelepiped $Q$ of
$\Lambda$.  Absolute integrability permits unfolding, and for
$\xi\in\Lambda^*$ the corresponding Fourier coefficient is
\[
 V^{-1}\int_Q F(z)e^{-2\pi i\langle\xi,z\rangle_{\mathbb R}}\,dz
       =V^{-1}\widehat\phi(\xi).
\]
The assumed absolute convergence implies that the right side of
\eqref{gt-Poisson-eq}, as a Fourier series in $z_0$, converges
uniformly to a continuous function with exactly those
coefficients.  The difference from $F$ has all Fourier
coefficients zero and therefore vanishes.  One elementary
justification of this last assertion is to identify
$\mathbb C^f/\Lambda$ with
$(\mathbb R/\mathbb Z)^{2f}$ using a lattice basis and convolve
with products of the one-dimensional Fej\'er kernels
\[
 F_N(t)=\frac1N\left|\sum_{j=0}^{N-1}e^{2\pi ijt}\right|^2.
\]
These kernels are nonnegative, have integral one, and their
mass outside any fixed neighborhood of $0$ tends to zero:
away from the integers the geometric-sum formula bounds
$F_N(t)$ by a constant divided by $N$.
Consequently their convolution with a continuous function on
the torus converges uniformly to that function.
All those convolutions are zero when all its Fourier
coefficients vanish.  This proves \eqref{gt-Poisson-eq}.
For real nonnegative $\phi$, the finite sum is at most the full
sum, and the triangle inequality bounds the nonzero Fourier
terms.  The zero coefficient is $\int\phi$, giving
\eqref{gt-Poisson-inequality}.
\end{proof}

\begin{corollary}\label{gt-actual-Poisson-error}
Let $P=P_{\nu,s}$, $c\ne0$, $\Lambda=c\sigma(P)$, and
$V=\operatorname{covol}\Lambda$.
Let $\phi$ be continuous, real, nonnegative, and satisfy
\eqref{gt-Poisson-decay}.
Suppose $B\geq0$ and nonnegative widths $x_v$ satisfy
\[
 V^{-1}\left|\widehat\phi
              \left(\frac2{gc}\overline{\sigma(\alpha)}\right)\right|
       \leq B W_{5,x}(\alpha)
       \quad(0\ne\alpha\in P^{\vee,g}),
 \qquad
 \prod_vx_v\leq(A/\Theta)^f(g\rho_E^{\nu+1})^{-f},
\]
where $A>0$ and $\Theta\geq A$.
For every $y_0\in E$ and finite subset $S\subseteq y_0+P$,
\begin{equation}\label{gt-actual-Poisson-error-eq}
 \sum_{y\in S}\phi(c\sigma(y))
       \leq V^{-1}\int\phi
                   +B\bigl(27(A/\Theta)^4\bigr)^f.
\end{equation}
The exact logarithmic-polynomial bound of
Corollary~\ref{gt-subgroup-sum} can replace the last factor.
If $E/F$ is unramified above $g$, then
\[
 V=|c|^{2f}2^{-f}g^{3f}\rho_E^{(2\nu+1)f}.
\]
In particular $c=\rho_E^{-\nu}$ gives
$V=(\rho_E/2)^fg^{3f}$.
\end{corollary}

\begin{proof}
The containing ideal for $P^{\vee,g}$ and its norm are exactly
\eqref{gt-dual-between} and \eqref{gt-dual-norm}.
Corollary~\ref{gt-subgroup-sum} therefore bounds the sum of
the nonzero Fourier coefficients divided by $V$.
Their sum is finite; the zero coefficient is finite because
$\phi$ is integrable.  Proposition~\ref{gt-Fourier-dual}
identifies all Fourier frequencies, with no omitted factor
of $2$, $g$, or $c$.  Theorem~\ref{gt-Poisson} now proves
\eqref{gt-actual-Poisson-error-eq}, and
Proposition~\ref{gt-exact-covolumes} gives the displayed volume.
\end{proof}

\subsection{Necessary hypotheses and the limits of this improvement}

\begin{example}\label{gt-critical-exponent}
The condition $K>2$ cannot be dropped from a uniform assertion of
convergence.  Take $E=\mathbb Q(i)$,
$\mathfrak a=\mathbb Z[i]$, and $x=1/\Theta<1$.  For $K=2$,
\[
 \sum_{0\ne\alpha\in\mathbb Z[i]}W_{2,x}(\alpha)
       =x^2\sum_{0\ne\alpha\in\mathbb Z[i]}|\alpha|^{-2}
       =\infty.
\]
Indeed, the $a$ points $a+bi$, $0\leq b<a$, have squared modulus
at most $2a^2$, so their contribution to the latter sum is at least
$1/(2a)$.  Summing over $a\geq1$ proves divergence.  The same
conclusion holds for $0<K\leq2$, since $0<q_x(\alpha)\leq1$.
The exponent in the actual theta application is $K=5$.
\end{example}

\begin{example}\label{gt-log-necessary}
The logarithmic factor at the full decay power cannot in general be
replaced by a constant independent of $\Theta$.  Take
$E=\mathbb Q(\sqrt2,i)$, so $f=2$, and
$\mathfrak a=\mathcal O_E$.  The algebraic integer
$u=1+\sqrt2$ is a unit, with inverse $\sqrt2-1$.
At the two complex places its absolute values are $u$ and $u^{-1}$.
Set $x_1=x_2=\Theta^{-1}$, where $\Theta>1$.
For every integer $n$ satisfying
$|n|\leq\log\Theta/\log u$, both absolute values of $u^n$
lie between $\Theta^{-1}$ and $\Theta$.  Hence
\[
 q_x(u^n)=\Theta^{-2},\qquad
 W_{K,x}(u^n)=\Theta^{-2K}.
\]
The powers are distinct, and therefore
\[
 \sum_{0\ne\alpha\in\mathcal O_E}W_{K,x}(\alpha)
 \geq
 \left(2\left\lfloor\frac{\log\Theta}{\log(1+\sqrt2)}
                         \right\rfloor+1\right)\Theta^{-2K}.
\]
This rules out a bound $C^2\Theta^{-2K}$ with a fixed $C$ for this
one field.  In degree four the degree-one logarithmic polynomial
in \eqref{gt-smoothed-exact} has the correct order of growth.
\end{example}

Two further restrictions are essential in the elementary count.
The zero element is excluded: for $\mathfrak a=\mathbb Z[i]$
and $0<R<1/3$, including zero would give at least one point,
whereas the proposed right side $9R^2$ is smaller than one.
Also the product lower bound must apply to differences of points.
A bound only on each individual point would allow arbitrarily many
distinct points on the complex unit circle and would not imply a
packing estimate.  The actual summation sets lie in fractional
ideals, so both restrictions are satisfied.

The unramifiedness assumption for the exact indices also has
content, although it was unnecessary for the containing ideals.
For $E=\mathbb Q(i)$, $g=2$, $\nu=0$, and $s=0$, one has
\[
 P_{0,0}=\{2(a+bi):4bi\in4\mathbb Z[i]\}=2\mathbb Z[i].
\]
Thus $[2\mathbb Z[i]:P_{0,0}]=1$, whereas the formula
$g^f$ would give $2$.  This prime is ramified in
$\mathbb Q(i)/\mathbb Q$, so it is excluded by
Proposition~\ref{gt-exact-covolumes} and by the actual source
hypotheses.  The example also explains why a proof that
diagonalizes the involution by dividing by $2$ would not
justify the asserted range of that proposition.

The source's stated smoothed-sum estimate has the weaker form
$(C\Theta^{-(K-5/2)})^f$; its proof already obtains an intermediate
bound with exponent $K-1$ before weakening it.  Thus retaining that
intermediate exponent is an improvement of the stated result, not
a correction of a false theorem.  The direct integration above
also gives the explicit logarithmic polynomial at exponent $K$,
smaller constants, and the exact convergence range $K>2$.
The actual application has $K=5$ and the width factor $2$ recorded
in Corollary~\ref{gt-actual-error}.

This proves the ideal-summation component and the global Poisson
step of the scalar orbital estimate, including its actual period
lattices, covolumes, trace and Fourier normalizations, and
discriminant exponents.  It does not by itself prove the width
inequalities \eqref{gt-dual-width}--\eqref{gt-lift-width}: those
come from the chosen archimedean test functions and their parameter
bounds.  The Fourier-coefficient bound for those particular
test functions must also be established.  Nor does the section
supply the finite-place mass identities or the remaining
larger-corank, nonscalar, and nongeneric orbital estimates.  The
degenerate diagonal estimate and the comparison with the ambient
automorphic dimension remain separate obligations.  No
counterexample to the actual ideal-sum or lattice-containment
statements has been found.


